\chapter{Decisiones de Diseño}
 
\section{Paradigmas y estrategia general}
 
Para el diseño de los algoritmos paralelos se adoptaron dos paradigmas de forma
complementaria:
 
\begin{itemize}
  \item \textbf{Paradigma de control} para las tareas de posterizado y borroneado: dado que
        ambas operan de forma independiente sobre la imagen de entrada (sin dependencias de
        datos entre sí), pueden ejecutarse en paralelo. Se decidió paralelizarlas asignando
        recursos de cómputo de forma independiente a cada una.
 
  \item \textbf{Paradigma de datos} para la distribución de los datos de entrada: dentro de
        cada tarea, los datos (filas o bloques de la imagen) se reparten entre los workers
        disponibles.
\end{itemize}
 
Dentro de la sección de borroneado, las dos subtareas (blur y highlight+umbralado) sí
presentan dependencia secuencial entre sí, por lo que se implementa un \textbf{pipeline de
datos}: primero se ejecuta la función de blur sobre toda la imagen y, sobre el resultado,
se aplica la función de highlight junto con la operación de umbralado.
 
\section{Algoritmo secuencial}
 
En todas las versiones del algoritmo se aplicó la técnica de \textbf{padding}: se añade un
borde auxiliar de píxeles con valor neutro (negro) alrededor de la matriz de entrada. Esto
elimina la necesidad de verificar en cada iteración si las celdas vecinas son válidas al
calcular el blur y el highlight, simplificando el código y mejorando la localidad de memoria.
 
\section{Algoritmo con memoria compartida (OpenMP)}
 
Se aplica descomposición de datos de entrada para cada tarea dentro del proceso.
 
La paralelización de las tareas de borroneado y posterizado se realiza con la directiva
 
\begin{verbatim}
#pragma omp parallel sections num_threads(2)
\end{verbatim}
 
que identifica las dos secciones a ejecutar en paralelo, combinada con
\texttt{\#pragma omp section} para cada una de ellas.
 
\textbf{Suma y posterizado.} Se usa \texttt{\#pragma omp parallel for schedule(static)}.
La planificación estática se eligió porque el costo computacional de cada iteración es
uniforme, lo que permite dividir el espacio de trabajo de forma equitativa sin sobrecarga
de administración dinámica.
 
\textbf{Blur y highlight.} La inicialización de datos se paraleliza con
\texttt{\#pragma omp parallel for schedule(static)}. El procesamiento principal emplea
 
\begin{verbatim}
#pragma omp parallel for collapse(2) schedule(guided)
\end{verbatim}
 
La cláusula \texttt{collapse(2)} fusiona los dos loops anidados en un único espacio de
iteraciones. La planificación \texttt{guided} se adopta porque el trabajo por iteración puede
ser variable o impredecible, y resulta más eficiente que \texttt{dynamic} en la administración
del balance de carga.
 
\section{Algoritmo con memoria distribuida (MPI)}
 
El proceso de rango 0 actúa como \textbf{orquestador}: recibe la imagen original, la
distribuye entre los workers y recolecta los resultados parciales. Los procesos worker se
dividen en dos grupos según la sección asignada: la mitad procesa el posterizado y la otra
mitad el borroneado.
 
Durante el análisis experimental se observó que paralelizar la \textbf{suma final} de las
matrices de posterizado y borroneado incrementaba el tiempo de ejecución de forma
contraproducente. Este comportamiento se atribuye a la latencia y overhead de comunicación
entre procesos, que supera el costo de la suma en sí misma. Por esta razón se decidió
\textbf{no paralelizar la suma final}, dejándola a cargo del proceso orquestador.
 
\section{Algoritmo híbrido (MPI + OpenMP)}
 
% TODO: completar con las decisiones de diseño del algoritmo híbrido
 